\documentclass[12pt]{article}
\usepackage[utf8]{inputenc}
\usepackage[T1]{fontenc}
\usepackage[margin=1in]{geometry}
\usepackage{setspace}
\usepackage{xcolor}
\usepackage[colorlinks=true, urlcolor=blue!60!black, linkcolor=blue!60!black]{hyperref}
\setlength{\parindent}{0em}
\setlength{\parskip}{0.45em}
\renewcommand{\baselinestretch}{1.0}
\begin{document}
\begin{center}
{\Large \textit{Baker v. General Motors Corp.}, 522 U.S. 222 (1998)\par}
\vspace{0.3em}
{\normalsize Case Brief --- One-Page Edition\par}
\end{center}

\textbf{Memory Jogger.} The gagged-witness case: judgments get exacting full faith and credit, but cannot control other courts' witnesses (\href{https://supreme.justia.com/cases/federal/us/522/222/}{522 U.S. 222}).

\textbf{Facts.} A Michigan consent injunction barred former GM engineer Ronald Elwell from testifying against GM. The sons of Beverly Garner, killed in a Missouri crash and fire in a Chevrolet Blazer, sued GM and subpoenaed Elwell.

\textbf{Procedural History.} The district court let Elwell testify and the jury awarded \$11.3 million; the Eighth Circuit reversed.

\textbf{Votes.} Reversed, 9--0 in the judgment; Ginsburg wrote for five, and Scalia and Kennedy concurred in the judgment.

\textbf{Holding.} Final judgments receive exacting full faith and credit with no public-policy exception, but enforcement measures follow forum law, and a decree cannot govern testimony in litigation by strangers to it.

\textbf{Reasoning.} Laws may yield to forum policy; judgments may not. But Michigan had no authority over the Bakers or Missouri's subpoenas, and GM sought to use Missouri as a surrogate tribunal. Scalia said judgments cannot be executed out of state; Kennedy looked only to Michigan's own preclusion law.

\textbf{Examples.} \emph{Same-Side:} subpoenaing a scientist gagged by another State's decree. \emph{Opposite-Side:} enforcing a sister-state money judgment despite contrary policy. \emph{Fence-Sitter:} judgments against shield-law physicians.

\textbf{Critique.} Kennedy warned that the injunction exception smuggles in a public-policy exception under a jurisdictional label.

\textbf{Key Quotation.} ``Regarding judgments, however, the full faith and credit obligation is exacting. A final judgment in one State, if rendered by a court with adjudicatory authority over the subject matter and persons governed by the judgment, qualifies for recognition throughout the land. \ldots But our decisions support no roving `public policy exception' to the full faith and credit due judgments.'' --- Justice Ginsburg, Opinion of the Court, 522 U.S. at 233.

\end{document}
